\section{A ten-positive-variable equation of degree 246}\label{sec:center246}

A second native companion improves degree rather than operation count.
It has ten positive unknowns, exact degree \textbf{246}, and a completely
paid schedule of \textbf{121 operations} (73 multiplications, 25 additions,
and 23 subtractions).  Its parent sharp-center equation has degree 254
and 97 operations \cite{gapfusion}.  This is a degree/cost tradeoff.
The following construction uses the same complete compact44 tuple as
Section~\ref{sec:positive}, but restores the independent coordinate $Y$:
\emph{the substitution $Y=q^3s$ is not made here}.

\subsection{The full polynomial and its two correction terms}

Let $J,F,\alpha,i,T,x,Y,n,\eta,\zeta\in\Pos$.  Define
\begin{align*}
 q&=\beta J+1,&X&=q+x,&g&=q(q-1),\\
 C_d&=q-F-\alpha,&B&=(qK_0+X)C_d-qF,\\
 \mu&=q^2+1,&t&=\mu Y,&L&=X+1,\\
 A&=tL+1,&\Delta&=A^2-1,&H&=4A-5,\\
 k&=\eta+\zeta,&c&=kY+\eta,&v&=Xt^2,\\
 D&=1+v(v+1)k^2,&E&=ic^2,\\
 F_s&=1+\Delta DE^2,&Q&=\Delta^2DE^2,&\delta&=Q-1.
\end{align*}
The remaining unchanged kernel and carrier data are
\begin{align*}
 R&=[(q-1)(q-F)+\alpha](q^2-1)+(M_C+qM_F)J,\\
 e&=k-R-1,&b&=c(T+\alpha)-R,&V&=F_sb-c,\\
 D_*&=QV^2-1,&P&=q^4(q-1),&M&=P\mu^2H,\\
 K_+&=X(3c-2X)+2,&W&=1+gb.
\end{align*}
The new polynomial abbreviations are
\begin{align}
 Q_1&=k^2\bigl[t^2LX^2(tL-1)+t(L^3-2L+1)-L^2+2\bigr],\label{eq:center-Q1}\\
 Q_3&=k^2(tX^2-3L+6),&\lambda&=8L^3,\label{eq:center-Q3}\\
 \Gamma&=E(8L^3AD-4Q_1-Q_3),&a&=\lambda^2DF_s.\label{eq:center-Gamma}
\end{align}
Set
\begin{equation}\label{eq:center-Z}
 \begin{aligned}
 Z={}&P\mu^2\bigl[H(n+V+\Gamma-bB)+K_+\bigr]\\
     &+H\bigl[\mu^2(gY-q^3B-(q-1)X)-Pe\bigr].
 \end{aligned}
\end{equation}
The complete equation is
\begin{equation}\label{eq:center246}
 \boxed{\mathcal N_{246}
 =\bigl(\delta Z^2-M^2D_*-\delta M^2a\bigr)^2
       -4\delta M^4aD_*=0.}
\end{equation}
All abbreviations are integer polynomials in the ten listed unknowns.
The corresponding literal graph and coefficient payloads are in
\pathfile{center246/}.  No additional quantified coordinate is hidden
in~\eqref{eq:center246}.

The construction comes from two exact polynomial divisions.  With
$\Omega=L^3D$ and $X=L-1$, expansion gives
\begin{align}
 \Omega&=AQ_1+r_1=A^3Q_3+r_3,\label{eq:center-divisions}\\
 r_1&=L^3+k^2(L^2-t-2),\nonumber\\
 r_3&=L^3+k^2\bigl[L^2t^2(L^2+5L-3t-12)
             +Lt(8L-3t-16)+3L-t-6\bigr].\nonumber
\end{align}
Thus $8\Omega A-4Q_1-Q_3$ retains the polynomial portions of the
first three terms of
$8\Omega\sqrt{A^2-1}=8\Omega A-4\Omega/A-\Omega/A^3-\cdots$.
The series is a discovery guide only: the construction and proof use
the exact identities~\eqref{eq:center-divisions}.
For a shorter evaluation schedule, put $p=A-1$, $u=tX^2$ and use
\begin{equation}\label{eq:center-paid}
 4Q_1+Q_3=k^2\bigl[u\{4p(p-1)+4L+9\}
                    +4(tX-L^2)-3L+14\bigr].
\end{equation}

\subsection{Positive bounds, root recovery, and transport}

\begin{lemma}\label{lem:center-bounds}
On the whole positive input domain, $0<\Gamma<Q$.
At every positive zero of~\eqref{eq:center246}, the recovered Pell roots
satisfy $\lambda\tau f+X/q^4<2Q<b/4$.
\end{lemma}
\begin{proof}
Here $L\ge33$ and $t\ge1025$.  Both $Q_1,Q_3$ are positive:
the last part of $Q_1/k^2$ is at least
$L^3-L^2-2L+3>0$, while
$Q_3/k^2\ge(L-2)(L-3)+1>0$.
The first term of $Q_1/k^2$ is at least $24t^3>t+2$, so
$r_1>-k^2(t+2)>-Q_1$ and~\eqref{eq:center-divisions} gives
$\Omega>(A-1)Q_1>Q_1$.  Also $Q_3<k^2tX^2<D\le\Omega$.
It follows that
\[
 0<\Gamma<\lambda ADE,\qquad
 \lambda A=8L^3(tL+1)\le16tL^4<t^4L^4<\Delta^2.
\]
Since $E\ge1$, this proves $\Gamma<Q$.

For completeness, root recovery does not require the desired transport
sign.  Since $F_s>c+1$ and $b$ is integral, $|V|>1$ and $D_*/\delta>0$.
The positive radicands $a$ and $D_*/\delta$ are unequal: equality
would give $\delta\lambda^2DF_s=QV^2-1$, impossible modulo $D>1$.
Writing $z=Z/M$, the norm therefore has $z\ne0$ and gives the rational
square root $(z^2+a-D_*/\delta)/(2z)$ of $a$.
As $\lambda>0$ and $\gcd(D,F_s)=1$, it follows that
$D=\tau^2$ and $F_s=f^2$ for positive integers $\tau,f$.
With $I=\tau E$, the unchanged adjacent-Pell descent gives an integer
$y>0$ and
\[
 f^2-\Delta I^2=1,\qquad Q=(\Delta I)^2,\qquad
 QV^2-(Q-1)y^2=1,
 \qquad z=\varepsilon\lambda\tau f+\nu y.
\]
Here $\varepsilon,\nu\in\{1,-1\}$.
The parent's pre-transport folding, first-index, and fifth-index
arguments now give $b>8Q$, $K_+>0$,
$0<y-V<b/4$, and $e/\mu^2<b/4$ \cite{gapfusion}.
None uses the new integral center or radical weight.
Finally $f<AI$ implies $\lambda\tau f<\lambda ADE<Q$;
also $X/q^4<Q$.  Their sum is less than $2Q<b/4$.
\end{proof}

The exact normalized carrier is
\begin{equation}\label{eq:center-carrier}
 \frac ZM=n+V+\Gamma+\frac{K_+}{H}+\frac Y{q^3}
          -\frac{WB}{g}-\frac X{q^4}-\frac e{\mu^2}.
\end{equation}
Combining it with the recovered signed roots and
Lemma~\ref{lem:center-bounds} gives
\[
 n<\frac{3b}{4}+\frac{WB}{g}.
\]
If $B\le-1$, then $WB/g\le-b-1/g$, contradicting $n>0$.
Thus $B\ge0$ and $C_d>0$.  The remaining ordered rank, parity,
coprimality, whole-product projection, and unscaled source reconstruction
are those of the pinned gap-fusion proof.  The center and weighted root
are integral; the fractional terms in~\eqref{eq:center-carrier} and
$W\equiv1\pmod g$ are unchanged.  Hence the same source61/native/controller
calibration recovers a finite ZFC refutation.  This is the inherited
native branch, not the direct C222 route of the principal quartic.

Conversely take the parent's fresh canonical kernel witnesses from a
refutation.  They supply integral $t_0=B/g\ge0$, $\rho=e/\mu^2\ge1$,
$s_Y=Y/q^3>0$, $\kappa=K_+/H>0$, $X/q^4>0$, and
$u=y-V>2Q$, with $s_Y,\kappa<c$.  Define the remaining coordinate by
\begin{equation}\label{eq:center-complete}
 n=Wt_0+\rho+\lambda\tau f-\Gamma+u-s_Y-\kappa+X/q^4.
\end{equation}
It is an integer and is positive because
$n>2Q-Q-2c>0$; here $Q>c^4>2c$ follows from $E=ic^2$ and $c\ge3$.
The carrier becomes $\lambda\tau f+y$, so the norm vanishes.
This proves existence equivalence by a fresh positive construction,
without claiming a positive pointwise inverse for every parent zero.

\subsection{Exact degree and a noncancelling coefficient}

Write $C_*:=\Gamma^2-a$.  Assign weights $1,3,1,5$ to the independent
polynomial variables $L,t,k,E$.  Then
$\degr A=4$, $\degr D=16$, $\degr\Omega=19$,
$\degr Q_1=15$, $\degr Q_3=7$, $\degr r_1=5$, and $\degr r_3=13$.
An exact identity derived from~\eqref{eq:center-divisions} is
\begin{equation}\label{eq:center-defect}
 \begin{aligned}
 A^2(C_*+64L^6D)=E^2\bigl[&64A^2\Omega r_1
       +16\Omega(r_3-2r_1)+16r_1^2\\
       &+8A^2Q_1Q_3+A^2Q_3^2\bigr].
 \end{aligned}
\end{equation}
The two degree-42 terms on the right have combined leading part
$-112L^7t^7k^4E^2$.  All others have smaller degree.
Since $A_0=Lt$, this proves
\begin{equation}\label{eq:center-defect-leading}
 \degr C_*=34,\qquad (C_*)_0=-112L_0^5t_0^5k_0^4E_0^2.
\end{equation}
The supplied certificate independently lists and checks all 52
coefficients of $C_*$.

In the actual ten inputs, the unchanged kernel has
\[
 \degr\delta=42,\quad\degr V=38,\quad\degr M=13,
 \quad\degr a=56,\quad\degr\Gamma=28.
\]
Write $Z=M(V+\Gamma)+Z_0$, where $\degr Z_0\le19$ and
$(Z_0)_0=-M_0b_0B_0$.  Put
\[
 \Theta=\delta M^2C_*-M^2(V^2-1)
       +2\delta M(V+\Gamma)Z_0+\delta Z_0^2.
\]
Direct expansion gives
\begin{equation}\label{eq:center-norm-expansion}
 \begin{aligned}
 \mathcal N_{246}={}&4\delta^2M^4V^2C_*
       +4\delta M^2V\Gamma\Theta+\Theta^2\\
       &-4\delta M^4aV^2+4\delta M^4a.
 \end{aligned}
\end{equation}
Here $\degr\Theta=112$, so the five displayed terms have degrees at most
$246,246,224,226,150$.  The first two terms both contribute at degree 246;
their possible cancellation must be checked.

Let $q_0=\beta J$, $X_0=q_0+x$, $t_0=q_0^2Y$, $k_0=\eta+\zeta$,
$E_0=ik_0^2Y^2$, $b_0=-q_0^3(q_0-F)$, and
$B_0=(q_0K_0+X_0)(q_0-F-\alpha)-q_0F$.
The complete leading form is
\[
 4\delta_0^2M_0^4V_0^2
 \bigl((C_*)_0-2\Gamma_0b_0B_0\bigr),
 \qquad\Gamma_0=8X_0^6t_0^5k_0^2E_0.
\]
Its first summand has degree ten in $i$ and equals
\[
 -7\cdot2^{14}q_0^{116}X_0^{29}i^{10}Y^{57}
                     k_0^{32}(q_0-F)^2;
\]
the second has degree nine in $i$.  Consequently
\begin{equation}\label{eq:center-coefficient}
 \boxed{[J^{116}x^{29}i^{10}Y^{57}\eta^{32}F^2]\,
          \mathcal N_{246}=-7\cdot2^{14}\beta^{116}\ne0.}
\end{equation}
This proves exact degree 246 for the actual fixed tuple, rather than
only for generic ports.

The companion checker passes 151 checks, including exact local polynomial
identities, complete coefficient and source pins, circuit liveness,
independent formula evaluations, and degree/leading-form transcription
checks over three finite fields.  Polynomial-part radical centering
already appears elsewhere in ProveIt, for example the unrestricted-integer
construction in \pathfile{87-direct-centered-round12/norm/PROOF.md}.
The new result is this concrete application and degree/cost point relative
to the frozen snapshot.  It inherits the parent's ordinary calibration;
no new end-to-end proof-assistant certificate or global minimality claim
is made.
